\section{Result and comparison}

The principal object of this report is the fixed ordinary polynomial
\begin{equation}\label{eq:main}
 \boxed{P_{238}(x_0,\ldots,x_{237})
       =\sum_{i=0}^{173} R_i(x_0,\ldots,x_{237})^2=0,
       \qquad (x_0,\ldots,x_{237})\in\mathbb Z^{238}.}
\end{equation}
All 174 quadratic residuals are printed in Appendix~\ref{app:residuals}.
The accompanying file \nolinkurl{equations/P238/equation.json} gives the
same complete list, and \nolinkurl{expanded_polynomial.json} in that
directory gives every collected coefficient. Thus the sum in
\eqref{eq:main} specifies one polynomial with fixed integer coefficients;
it is not a variable-length system or a universal-polynomial placeholder.

\begin{theorem}[Consistency equivalence]\label{thm:main}
For the explicitly specified polynomial \eqref{eq:main}, the ordinary
arithmetic proof below, composed with the identified inherited numerical
calculus calibration, establishes
\[
 \mathrm{ZFC}\vdash
 \left(\neg\exists\boldsymbol x\in\mathbb Z^{238}\;
                 P_{238}(\boldsymbol x)=0\right)
 \longleftrightarrow\Con(\mathrm{ZFC}).
\]
The same assertion holds for the two accompanying 239-variable
polynomials, with their specified domains and coefficients.
\end{theorem}

Here consistency is syntactic consistency: no finite proof of falsity from
the fixed ZFC axiom predicate. The proof assumes neither consistency nor
inconsistency. It supplies an ordinary mathematical argument formalizable
in ZFC; it does not supply a newly compiled end-to-end Lean/Rocq theorem
or an encoded object-language ZFC proof string.

\subsection{Exactly what improved}

\begin{center}\small
\begin{tabular}{lrrrr}
\toprule
Equation over $\mathbb Z$ & Variables & Rows & Monomials & Height\\
\midrule
Incoming Shared Factorial Error &240&177&21,482&6,144\\
Incoming factorial-free, sparse &242&178&5,018&7,968\\
Incoming factorial-free, dense &241&177&12,774&7,968\\
\midrule
New $P_{238}$ &\textbf{238}&\textbf{174}&\textbf{13,149}&6,144\\
New $P_{239}^{\rm sparse}$ &239&175&\textbf{4,963}&7,968\\
New $P_{239}^{\rm low}$ &239&175&6,723&\textbf{4,432}\\
\bottomrule
\end{tabular}
\end{center}

Every polynomial in the table has exact degree four. ``Rows'' means the
number of residuals whose squares are summed; ``monomials'' counts terms
after complete collection; height is the largest absolute collected
coefficient. Every coefficient of an individual new residual has absolute
value at most 64. Every declared variable occurs in the collected polynomial.
The comparison equations are retained and recollected by the verifier.

Against the smallest-dimension direct quartic in the seven incoming
archives at the pinned snapshot, $P_{238}$ removes two variables, three
rows, and 8,333 monomials, with unchanged height. The new low-height
239-variable equation improves all four size measures of that 240-variable
comparison. Against its immediate factorial-free sparse predecessor,
$P_{239}^{\rm sparse}$ removes three variables, three rows and 55 monomials,
with unchanged height.

These are improvements to the direct numerical-calculus construction.
The repository also has much smaller variable counts in other coefficient
and degree architectures, including a 37-integer quartic and high-degree
few-variable forms~\cite{repo}. The present result does not establish a
global simplicity minimum. Its advantages are a fully specified quartic,
small literal coefficients, and a direct proof-calculus interpretation.

\subsection{The mathematical contribution}

The immediate parent has already removed factorial authenticity from its
CRT construction~\cite{factorialfree}. This report removes three more
arithmetic coordinates by a coordinated argument:

\begin{enumerate}
\item The single growth value $F=t^t$ suffices both to distinguish bounded
      proof-label products and to serve as the CRT radius. Thus $R=F$
      replaces $R=F^2$.
\item An even CRT index and a universally nonunit affine majorant make the
      power base nonzero before any representative sign is known. This
      removes the separate square restriction on the radix.
\item An existing power request recovers the sign of the beta spacing.
      Its odd multiplier can therefore be affine, instead of squared.
\end{enumerate}

The elementary LCM bound and polynomial coefficient arguments are proved
here, without a priority claim for those general facts. The new contribution
is their use together with a noncircular sign-recovery order in this explicit
source. The final row-basis reductions are reversible integer linear algebra,
not additional logical assumptions.

\section{Construction and inherited interfaces}\label{sec:interfaces}

The source uses the fixed numerical proof calculus $C_{222}$, with truth
code 0 and falsity code 1. A finite multiset of locally valid, single-conclusion
rules is certified arithmetically. Beta coding stores the finite accumulator
sequence; a binomial CRT certificate checks all local rows at once; a shared
Pell construction certifies the finite set of required powers. This route
uses neither a simulation of K215 nor a specialization of Jones's universal
equation. General Diophantine representation results provide historical
context, rather than an unspecified coefficient in this construction~\cite{jones}.

The unchanged source supplies fourteen principal Pell pairs, eighteen
power requests, a signed frame, a semantic anchor, and one common norm
bound for eleven modular remainders. The following interfaces are retained
with their ordinary proofs and source identities:

\begin{center}\small
\begin{tabular}{p{0.28\linewidth}p{0.64\linewidth}}
\toprule
Interface & Consequence used here\\
\midrule
Signed Pell frame and anchor & Canonical signed-index Pell coordinates,
before arithmetic signs; every actual base, index and output has absolute
value below $W$, where $W\ge258$ and $|A|>W^W$.\\
Local constructor ledger & Derived-label bound $h_*\le H^{2^{22}}$ after
representatives are nonnegative; local nonnegative row value below
$(2+t+x+y+F)^{2^{23}}$.\\
Horn calibration & Valid local rules and the recovered label balance imply
a finite $C_{222}$ derivation of code 1.\\
Logical calibration & $C_{222}$ refutability is equivalent to syntactic ZF
refutability, and then to syntactic ZFC refutability.\\
\bottomrule
\end{tabular}
\end{center}

The named logical interfaces are
\begin{center}\small
\nolinkurl{KEndpoint.c222_iff_zf_refutation},\\
\nolinkurl{ConstructibleChoice.zf_refutation_iff_zfc_refutation}.
\end{center}
Their source files are included under the inherited calibration directory.
The first theorem is about the numerical calculus itself; the machine
simulation hypotheses in other theorems of its source file are not premises
of this theorem. The second interface uses the constructible-universe
interpretation to transport ZFC refutations to ZF refutations. This report
does not rebuild their complete transitive proof-assistant environment.

\section{The new arithmetic choices}\label{sec:choices}

Use unrestricted integer certificate coordinates and the semantic values
\begin{equation}\label{eq:choices}
\begin{gathered}
 t=16+4h^2,\qquad F=t^t,\qquad R=F,\qquad C=2c,\\
 H=2+4\left(C+\sum_{j=1}^{8}A_j\right),\qquad
 y=F(2v+1).
\end{gathered}
\end{equation}
The nine actual CRT representatives are $C,A_1,\ldots,A_8$.
Historical stored coordinates may be these values minus one. The literal
source mutation consistently changes the actual index to $C=2c$.

The retained endpoint and CRT equations include
\begin{align}
 Bh_B&=C+F+1,\label{eq:raw-index}\\
 x&=F+1+(y+1)u_0,\label{eq:first}\\
 x&=1+[1+(t+1)y]u_t.\label{eq:last}
\end{align}
The power graphs request
\[
 F=t^t,\qquad q=(BH)^x,\qquad N=B^x,\qquad
 M=q^x,\qquad D_y=q^{y+1},
\]
as well as the binomial numerator and denominator powers. At the start of
soundness, these displayed semantics must be recovered from their polynomial
certificates. In particular, no positive value of $B$, $q$, $x$, $y$ or a
representative is assumed.

Variable exponentiation is explanatory notation in this section. The final
residuals contain only integer constants, addition, subtraction and
multiplication. The existing first power slot is reused with base and
index both $t$; no nineteenth power request is introduced. The old affine
base slot is retained with value $t$, so the frame counts its square twice.
This harmless duplication strengthens the bound and omits no actual atom.

\section{A sharp enough growth value for label separation}\label{sec:labels}

For a label $j\in\{0,\ldots,t\}$ set $W_j=1+jF$, so $W_0=1$.
There are $t$ local rows, each with one conclusion and at most two premises.
The desired product balance is
\begin{equation}\label{eq:products}
 \prod_{i<t}W_{v_i}=(1+F)\prod_{i<t}W_{p_i}W_{q_i}.
\end{equation}
The symbol $q_i$ here denotes a premise label, not the arithmetic radix $q$.

\begin{lemma}[Label separation at $F=t^t$]\label{lem:labels}
For $t\ge1$ and $F=t^t$, equality \eqref{eq:products} implies that, for
each positive label $j$, its conclusion multiplicity minus its premise
multiplicity is 1 when $j=1$ and 0 otherwise.
\end{lemma}
\begin{proof}
Discard unit factors with label zero. If the two remaining factor counts
differ, let $r$ be their minimum. One side originally has at most $t$
factors, so $r\le t$. The side with more factors is at least
$(F+1)^{r+1}$; the other is at most
\[
 (1+tF)^r\le t^r(F+1)^r.
\]
This contradicts $F+1=t^t+1>t^r$. Thus both sides have the same number
$m\le t$ of nonunit factors.

Replace $F$ by an indeterminate $z$. For $m>0$ and $0\le k\le m$,
the coefficient of $z^k$ in either product of $m$ factors $1+jz$,
$1\le j\le t$, is a positive integer. It is at most
\[
 \binom{m}{k}t^k
 =\binom{m}{m-k}t^k
 \le m^{m-k}t^k\le t^m\le F.
\]
Because both coefficients are positive, the corresponding difference has
absolute value at most $t^m-1<F$. This strict inequality remains valid
at the boundary $m=t$; using only an upper bound by $F$ without positivity
would not suffice. If $m=0$, both products are already the polynomial 1.

A nonzero integer polynomial whose coefficients have absolute value below
$F$ cannot vanish at $F$: divide its evaluation by $F$ to the degree of
its lowest nonzero coefficient, then reduce modulo $F$. That coefficient
would be a nonzero multiple of $F$ of smaller absolute value. Therefore
the evaluated equality lifts to an identity in $\mathbb Z[z]$.
The distinct labels give distinct roots $-1/j$ of the linear factors
$1+jz$. Unique factorization over $\mathbb Q[z]$ gives the claimed
multiplicities.
\end{proof}

This replaces the parent's sufficient condition $F>(t+1)^t$. That older
condition is false for the new $F$ and is not silently retained. No
coprimality of the evaluated integers $W_j$ is needed.

\section{The CRT radius can be linear}\label{sec:radius}

Write $\Lambda_t=\operatorname{lcm}(1,\ldots,t)$.

\begin{lemma}\label{lem:lcm}
For every $t\ge1$, $\Lambda_t\le4^{t-1}$.
\end{lemma}
\begin{proof}
Let $m=\lceil t/2\rceil$. For each prime $p$, the exponent of $p$ in
$\Lambda_t$ exceeds that in $\Lambda_m$ by at most one. If it exceeds
it, a power $p^a$ lies above $m$ and at most $t$. Its contribution to
the factorial valuation of $\binom{t}{\lfloor t/2\rfloor}$ is one;
all other contributions are nonnegative. Consequently
\[
 \Lambda_t\mid\Lambda_m\binom{t}{\lfloor t/2\rfloor}.
\]
For $t=2n$, use strong induction and $\binom{2n}{n}\le4^n$ to obtain
$\Lambda_{2n}\le4^{n-1}4^n=4^{2n-1}$.
For $t=2n+1>1$, the two equal middle binomial coefficients have sum at
most $2^{2n+1}$, so each is at most $4^n$. Hence
$\Lambda_{2n+1}\le4^n4^n=4^{2n}$. The base case $t=1$ is immediate.
\end{proof}

It follows directly from $t\ge16$ that
\begin{equation}\label{eq:radius}
 R=F=t^t\ge16^t>16^{t-1}\ge\Lambda_t^2.
\end{equation}
Thus the existing growth value already meets the binomial-CRT radius
requirement. The separate square register $F^2$ is unnecessary.

For clarity, the two prime-power facts used by that compiler are reproduced
next. The classical binomial-LCM divisibility in their proof is also
related to Farhi's binomial-row identity~\cite{farhi}.

\begin{lemma}[Two-sided prime-power coverage]\label{lem:coverage}
Let $R\ge\Lambda_t^2$ and $B=\binom{R+t}{t}$. Then:
\begin{enumerate}
\item For every $i\in\{1,\ldots,t\}$, some full prime power
      $p^e\Vert B$ divides $R+i$ and satisfies $p^e>t$.
\item Every full prime power $p^e\Vert B$ divides $R+i$ for at least
      one $i\in\{1,\ldots,t\}$.
\end{enumerate}
\end{lemma}
\begin{proof}
First, $t\binom{t-1}{i-1}\mid\Lambda_t$. At a prime power $p^a$, the
factorial valuation summand
\[
 \left\lfloor\frac{t}{p^a}\right\rfloor
 -\left\lfloor\frac{i-1}{p^a}\right\rfloor
 -\left\lfloor\frac{t-i}{p^a}\right\rfloor
\]
is 0 or 1 and vanishes for $p^a>t$.

Fix $i$. Since $R+i>\Lambda_t^2$, there is a prime with
$v=v_p(R+i)>2a$, where $a=v_p(\Lambda_t)$. For $j\ne i$,
$v_p(R+j)=v_p(j-i)$, because $p^v>t>|j-i|$.
Therefore
\[
 e=v_p(B)=v-v_p\!\left(t\binom{t-1}{i-1}\right).
\]
The subtracted quantity lies between 0 and $a$, so $e>a$ and $e\le v$.
Thus $p^e>t$ and $p^e\mid R+i$, proving the first assertion.

For the second assertion, set $v=\max_i v_p(R+i)$. The valuation of
$B$ is a sum of 0-or-1 factorial summands. Above level $v$, the interval
$R+1,\ldots,R+t$ contains no multiple of that prime power; all such
summands vanish. Hence $e=v_p(B)\le v$. A row attaining the maximum
has $p^e\mid R+i$.
\end{proof}

The second conclusion is essential for completeness. Choosing just one
selector for each row would leave unspecified how to satisfy the CRT
congruences at the remaining primes of $B$.

\section{Signed power certification and the recovery order}\label{sec:bootstrap}

The retained Pell interface can be used with a negative base. Let
$X_k(A),Y_k(A)$ denote the usual Pell polynomials, continued to negative
indices by $X_{-k}=X_k$, $Y_{-k}=-Y_k$. For a requested base $n\ne0$,
put $T_n=2An-n^2-1$. The power congruence and Pell recurrence give
\[
 m\equiv n^k\pmod{T_n}\quad(k\ge0),\qquad
 mn^{-k}\equiv1\pmod{T_n}\quad(k<0).
\]
The signed frame bounds the relevant absolute errors by less than
$|T_n|$. For $|n|\ge2$, a negative $k$ would therefore force the
impossible equality $mn^{-k}=1$ in integers. A nonnegative $k$ forces
$m=n^k$. This is the precise inherited signed-power interface used below;
no assertion at base zero is needed.

\begin{proposition}[Acyclic sign recovery]\label{prop:bootstrap}
Every integer zero of the changed source satisfies
\[
 F>2^{32},\quad y\ge F>0,\quad x\ge2+(t+1)y,
 \quad C,A_j\ge0,\quad B>0,\quad q>0,
\]
with the intended power and binomial meanings.
\end{proposition}
\begin{proof}
The steps must be taken in the following order.

\textbf{1. Recover $F$ and exclude a zero base.}
The request with known positive base and index $t$ gives
$F=t^t\ge16^{16}=2^{64}$. Since $t=16+4h^2$, $F$ is even.
The right side of \eqref{eq:raw-index} is odd even if $C$ is negative.
Thus $B$ and $h_B$ are odd and nonzero. Independently,
$H\equiv2\pmod4$ for all signed representatives, so $|H|\ge2$.
The power base therefore satisfies $|BH|\ge2$ without a positivity
assumption on any representative.

\textbf{2. Recover $x$ and make the radix nonunit.}
The request at base $BH$ gives $x\ge0$ and $q=(BH)^x$.
The affine beta definition gives $|y|\ge F$, because $2v+1\ne0$.
Hence
\[
 |1+(t+1)y|\ge(t+1)F-1>1.
\]
If $x=0$, the last endpoint would make this integer divide $-1$,
which is impossible. Thus $x\ge1$ and $Q=|q|\ge2$.

\textbf{3. Recover the beta sign from an existing power.}
The already present request with base $q$ and exponent $y+1$ is now
usable. It forces $y+1\ge0$. A negative $y$ would satisfy
$y\le-F<-1$, so $y\ge F>0$. This is the step that permits the
unsquared odd beta multiplier.

\textbf{4. Obtain the large exponent before representative signs.}
If $u_0\le-1$, \eqref{eq:first} would give
$x\le F-y\le0$. Therefore $u_0\ge0$ and $x\ge F+1$.
Equation~\eqref{eq:last} then gives $u_t\ge1$ and
\begin{equation}\label{eq:xgrowth}
 x\ge2+(t+1)y\ge2+(t+1)F.
\end{equation}
Consequently $Q\ge2^x$ and $|q+1|\ge Q-1\ge2$.

\textbf{5. Recover the representative signs.}
Each actual representative is an exponent in a request at base $q+1$.
The signed-power lemma gives $C,A_1,\ldots,A_8\ge0$, and now $H\ge2$
is positive. The sign of $B$, and therefore possibly that of $q$, is
still undetermined at this point.

\textbf{6. Authenticate the first binomial digit in either radix sign.}
Set $n=F+t$ and $b_j=\binom{n}{j}$. For $0\le j\le t$,
\[
 b_j\le(F+t)^t\le(F^2+F)^t\le2^{t(F-1)}.
\]
The last elementary inequality holds for $F\ge7$. From
\eqref{eq:xgrowth},
\[
 \frac{2^{t(F-1)}}{Q}\le2^{-F-t-2}.
\]
Thus the signed low part
$r_* =\sum_{j<t}b_jq^j$ has $|r_*|<Q^t/32$, and $b_t<Q/32$.
The binomial theorem is valid also for $q<0$.

To see exactly what the common norm enforces, write a denominator
$D=q^k$, its ratio $a=M/D$, and a normalized remainder $\rho$.
The literal rows give $Da=M$ and $aU=MJ+\rho$. Since $a\ne0$,
$a\mid\rho$; putting $r=\rho/a$ gives $U=DJ+r$ with $r\in\mathbb Z$.
The common norm is
\begin{equation}\label{eq:common}
 M^2=1+4\sum_{i=1}^{11}\rho_i^2+\alpha^2+\beta^2+\gamma(\gamma+1).
\end{equation}
Its last three terms are nonnegative for every integer input. Therefore
$|\rho_i|<|M|/2$ and $|r_i|<|D_i|/2$, regardless of denominator signs.
For the first binomial numerator, $D=q^t$ divides $r-r_*$, while
$|r-r_*|<|D|/2+|D|/32<|D|$. Hence $r=r_*$.

The next digit satisfies $B\equiv b_t\pmod q$. Put $a=|BH|\ge2$.
As $x\ge3$, $|B|\le a<Q/2$. Hence
$|B-b_t|<Q/2+Q/32<Q$ and
\begin{equation}\label{eq:binomial}
 B=b_t=\binom{F+t}{t}>0.
\end{equation}
Now $q=(BH)^x>0$, $B\ge2$, and the $B$ request gives $N=B^x>0$.
This finishes the sign recovery without using the binomial conclusion
to establish a sign needed in its own proof.
\end{proof}

\section{Recovering all local rows and a refutation}\label{sec:sound}

The new positive $H$ dominates the old majorant $2+C+\sum A_j$.
For an ordinary representative $a\ge0$, every needed low coefficient
satisfies $\binom aj\le\max(1,a)^x\le H^x$. This includes zero
representatives and the constant coefficient. The inherited derived-label
bound is $h_*\le H^{\ell}$ with $\ell=2^{22}$. Since $F>\ell$ and
$x>(t+1)F$, its needed coefficients also lie below $H^x$.
Since $q=B^xH^x$ and $B\ge2$,
\[
 \frac{H^x}{q-1}<\frac{2}{2^x}<\frac1{32}.
\]
Thus all ten remaining normalized binomial remainders are identified
exactly by the same argument used for the first digit.

\subsection{The binomial CRT implication}

The elementary range fact is: if $p$ is prime and $p^e\mid\binom ak$, with $a\ge0$ and
$k,e\ge1$, then $a\bmod p^e<k$. If $a<k$ it is immediate. Otherwise
the factorial valuation is a sum of 0-or-1 terms. If the $k$ consecutive
factors $a-k+1,\ldots,a$ contain no multiple of $p^e$, all terms from
level $e$ onward vanish, leaving valuation at most $e-1$, a contradiction.

Use $N=B^x$. Fix a row $i\in\{1,\ldots,t\}$ and its selector
$p^e\Vert B$ from Lemma~\ref{lem:coverage}. The exact index range
$N\mid\binom Ct$ gives $C\equiv s\pmod{p^{xe}}$ with $0\le s<t$.
The raw congruence $B\mid C+F+1$, reduced modulo $p^e\mid F+i$,
gives $s\equiv i-1\pmod{p^e}$. Since $p^e>t$, $s=i-1$ exactly.
The other binomial ranges similarly recover the bounded local parameters.

The unchanged local polynomial $L$ is nonnegative and has the bound
\[
 L<(2+t+x+y+F)^E,\qquad E=2^{23}.
\]
By \eqref{eq:xgrowth} and $y\ge F$, the parenthesis is less than $2x$;
also $x>2^{32}$. The elementary dyadic estimate
$E\log_2(2x)<x$ therefore gives $L<2^x$. For example, if
$2^k\le x<2^{k+1}$ and $k\ge32$, then
$E\log_2(2x)<2^{23}(k+2)<2^k\le x$; the middle inequality holds at
$k=32$ and persists by induction. The local congruence makes $p^{xe}$
divide $L$, but $p^{xe}\ge2^x$, so $L=0$. Every required row is valid.
The derived label is a polynomial of the independent columns, so its
congruence and nonnegative size bound identify the same local label.

\subsection{The exceptional beta boundary and the Horn argument}

The affine choice permits $y=F$ algebraically. At that boundary, the
first endpoint has beta remainder 0 modulo $y+1$, while the last has
remainder 1. The decoded local updates have the form
\[
 S_{i+1}W_{v_i}=S_iW_{p_i}W_{q_i},
\]
with every weight positive. An initial zero propagates to zero, contradicting
the terminal value. Thus every actual zero has $y>F$, initial accumulator
$F+1$, and terminal accumulator 1. All accumulators are positive;
multiplication and cancellation yield \eqref{eq:products}.

Lemma~\ref{lem:labels} gives the exact multiplicity balance. Take the
finite closure of truth under the rule instances appearing in the certificate.
If falsity were outside that closure, each row with an outside conclusion
would have at least one outside premise. Summing conclusions minus premises
over outside labels would then be nonpositive. The recovered balance makes
that sum one, a contradiction. Thus there is a finite $C_{222}$ refutation.
The identified calibration gives a finite ZFC refutation. Cyclic rule
multisets do not invalidate this argument; the single-conclusion property
is the relevant restriction.

\section{Constructive completeness}\label{sec:complete}

Start with a finite $C_{222}$ refutation. Expand it to a finite proof tree
and choose $t=16+4h^2$ beyond its size, schema parameters and label bounds.
This cofinal choice is made before the large arithmetic witnesses. Set
$F=t^t$ and construct the finite pending-multiset accumulator path with
weights $1+jF$, padding with truth to exactly $t$ rows.

Let $O$ be the odd part of $t!$. Choose an arbitrarily large positive odd
multiple $w$ of $O$ and put
\[
 y=Fw,\qquad v=(w-1)/2.
\]
Since $v_2(F)=t\,v_2(t)\ge2t>v_2(t!)$, this $y$ is an even multiple
of $t!$. Increase $w$ so that $y>F$ and every accumulator is at most $y$.
The moduli $1+(i+1)y$, $0\le i\le t$, are pairwise coprime: a common
prime would divide a nonzero index difference at most $t$, hence divide
$y$, contrary to its divisibility into $1+(i+1)y$. Choose a sufficiently
large positive beta code $x$ in the required CRT class. Its local beta
quotients are nonnegative and strictly below $x$ because their denominators
are at least two. The genuine endpoints give \eqref{eq:xgrowth}.

Set $B=\binom{F+t}{t}$ and $N=B^x$. To construct the second CRT stage,
for \emph{every} prime $p\mid B$ choose a compatible row using the second
part of Lemma~\ref{lem:coverage}. Prescribe that row's local tuple and
index modulo $p^{2xe}$, where $e=v_p(B)$. The pairwise coprime moduli have
product $N^2$. Range binomials are divisible by $p^{xe}$: their falling
factorials contain the selected factor divisible by $p^{2xe}$, and division
by $k!$ loses at most $v_p(k!)\le k-1<x\le xe$. The zero-binomial case
when the upper argument is smaller than $k$ is immediate. This supplies
all required global congruences, including those at unused selectors.

The binomial $B$ is odd. Indeed $v_2(F)>\lfloor\log_2t\rfloor$, so
$v_2(F+i)=v_2(i)$ for $1\le i\le t$; subtracting $v_2(t!)$ gives zero.
Thus $N^2$ is odd, and adding $N^2$ changes a representative's parity
without changing any CRT congruence. Choose $C$ even, giving $c=C/2$.
As in the parent, choose the dominant schema representative $a$ odd in
$[2N^2,4N^2)$ and $b$ even; choose the other representatives positive and
at most $2N^2$. The inherited constructor satisfies
\[
 h_*=1+a+b\pmod2,\qquad h_*\ge2a^2.
\]
It therefore supplies the same even dominant semantic anchor.
The changed front index is smaller: $F+t\le\binom{F+t}{t}=B<N^2$.
All other requested indices are bounded as in the parent construction.

Only now define $H=2+4(C+\sum A_j)$ and $q=(BH)^x$. This order avoids
any circular choice of a CRT modulus. Here $q$ is even, which is allowed;
the retained range argument uses $N\mid q$, since $q=NH^x$, and does
not require odd $q$. All four binomial denominator groups divide $M=q^x$.
Their genuine normalized remainders satisfy $0\le\rho_i<M/32$, so
\[
 4\sum_{i=1}^{11}\rho_i^2<\frac{44}{1024}M^2<M^2.
\]
The remaining integer gap in \eqref{eq:common} is nonnegative and is
representable as $a^2+b^2+c(c+1)$. This standard consequence of the
three-square theorem follows by representing $4g+1$: exactly one square
is odd and two are even, and division by four gives the displayed form.

The retained constructive Pell frame and semantic anchor now supply the
finite family of genuine power certificates. The first request still has
the existing index $t$. The frame's duplicated base/index atom is harmless.
This constructs a zero of the changed source. Exact finite transformations
then give zeros of each released polynomial.

\section{From the source to the literal quartics}\label{sec:finite}

The complete source mutation changes eleven surviving quadratic rows and
deletes three rows. All eighteen power requests remain. The deleted
registers are the squared radius, the private radix square root, and the
beta multiplier's square wire. Historical unused slots remain in the raw
array only to make the source comparison exact; they are absent from the
final coordinate list.

The finite construction then uses the retained Horn factoring identity,
202 unit-monic eliminations, deletion of 117 unused slots, and 69 integral
translations to obtain the 242-coordinate intermediate system with
178 quadratic rows. Six credited integer coordinate shears and exact row
identities expose three affine eliminations for the sparse 239-variable
system, or four for the 238-variable system. The last step is a recorded
sequence of row replacements $R_i\leftarrow R_i+kR_j$, $i\ne j$,
$k\in\mathbb Z$. These have the explicit inverse
$R_i\leftarrow R_i-kR_j$ and preserve the common zero set. They change
the sum-of-squares polynomial and can reduce its coefficient height.

For every unit-monic equation $\varepsilon z+f=0$, $\varepsilon=\pm1$,
the removed coordinate is restored as the integer polynomial
$z=-\varepsilon f$. The certificates are acyclic and the checker verifies
that every retained row stays quadratic. The final affine transformations
give a polynomial bijection of zero sets with the intermediate system.
Deleting unused coordinates earlier gives existence equivalence with free
fibers; it is not described as a bijection.

No expansion estimate is used for the reported metrics. Every final
coefficient is collected exactly. A nonzero degree-four term establishes
exact degree four, and every declared coordinate is checked to occur.

\section{Verification, reproducibility and proof boundary}\label{sec:verification}

The archive separates three kinds of evidence.

\paragraph{Unbounded mathematics.}
The proofs in this report and the two independent derivation/review notes
cover label separation, the LCM/radius inequality, the ordered signed
bootstrap, binomial decoding, both completeness constructions, and the
composition with the inherited logical interfaces. Those arguments, rather
than a finite search, establish the quantified equivalence.

\paragraph{Exact finite algebra.}
A separately written checker imports no producer or optimizer. It replays
the entire new source delta; seventeen affine-square identities and eighteen
Horn rewrites with private-wire checks; every monic restoration and integral
translation; and the complete final coordinate/row certificates. For each
endpoint it pulls back all 178 intermediate residuals through the inverse
row operations and compares every polynomial coefficient. Both ordinary
and optimized Python execution pass. Targeted controls reject a weakened
$C=2c$ chart, a changed Horn coefficient despite an updated hash, and a
monic restoration changed by one.

\paragraph{Supplementary arithmetic diagnostics.}
Independent finite checks cover the LCM recurrences through 500, 168
binomial-CRT examples, 756 row selectors, 1,269 full-prime support checks,
17,576 label-multiset coefficient/injectivity cases, and 4,878 unequal-factor
count comparisons. They also check actual parity/beta completions and a
counterexample to removing the CRT radius hypothesis. These diagnostics
help catch transcription mistakes; they are not substitutes for the proofs.

Run, with Python 3.10 or newer and its standard library,
\begin{verbatim}
python3 -B verify.py
python3 -B verify.py --rebuild
python3 -B -O verify.py
\end{verbatim}
The rebuild mode reconstructs the source and all selected coefficient files
byte for byte in a temporary directory. It replays frozen row certificates;
it does not rerun an optimizer or search for a refutation. To evaluate a
supplied 238-integer tuple stored as a JSON list, run
\begin{verbatim}
python3 -B evaluate.py point.json
\end{verbatim}
The evaluator is exact and makes no attempt to decide solvability.

Rebuild the standalone article with \texttt{python3 -B build\_report.py}.
The report source includes the complete printed polynomial definition.
Its generated portions are compared with the authoritative JSON. Different
TeX installations can change PDF metadata and bytes; the mathematical files
and deterministic source reconstruction are the reproducibility targets.

\section{Assessment and further work}

The achieved result is a smaller direct quartic with a more economical
arithmetic justification: one growth parameter performs two jobs, and
existing powers enforce signs that previously required separate squares.
The sharp point in the label proof is the positivity of both coefficient
lists, which allows $F=t^t$ exactly. The sharp point in the bootstrap is
using $q^{y+1}$ before using the representative powers at base $q+1$.
Interchanging those steps would reintroduce a circular argument.

This is a concrete new reduction relative to the inspected source lineage.
It is not a demonstrated breakthrough in the global minimum number of
Diophantine unknowns. The most substantial remaining arithmetic cost is
the shared certification of eighteen powers; a different simultaneous-power
certificate could have a much larger effect than another coordinate or
row-basis adjustment. The compact new lemmas and the exact source delta
also provide clear targets for future formalization. No unproved shortcut
is included in the released endpoints.

\begingroup\raggedright
\begin{thebibliography}{9}
\bibitem{repo} V. Reshetnikov and ProveIt contributors,
\emph{ProveIt}, comparison snapshot
\texttt{7a8da414376a80b90ecf799f37ceb26fb9ac674e}, 25 October 2026.
\url{https://github.com/VladimirReshetnikov/ProveIt/tree/7a8da414376a80b90ecf799f37ceb26fb9ac674e/Computability/DiophantineConsistency/con-zfc-diophantine-forms}.
\bibitem{factorialfree} ProveIt incoming research,
\emph{Removing factorial authenticity: binomial prime-power certificates
for Diophantine consistency},
\nolinkurl{ConZFC_Factorial_Free_Binomial_CRT_2026-10-25.zip},
retained in \nolinkurl{inherited/factorial_free/}.
\url{https://github.com/VladimirReshetnikov/ProveIt/blob/7a8da414376a80b90ecf799f37ceb26fb9ac674e/docs/incoming/ConZFC_Factorial_Free_Binomial_CRT_2026-10-25.zip}.
\bibitem{shared} ProveIt incoming research,
\emph{Shared factorial errors: explicit Diophantine equations for Con(ZFC)},
\nolinkurl{ConZFC_Shared_Factorial_Error_2026-10-25.zip}.
\url{https://github.com/VladimirReshetnikov/ProveIt/blob/7a8da414376a80b90ecf799f37ceb26fb9ac674e/docs/incoming/ConZFC_Shared_Factorial_Error_2026-10-25.zip}.
\bibitem{farhi} B. Farhi,
\emph{An identity involving the least common multiple of binomial
coefficients and its application}, 2009, arXiv:0906.2295v2.
\url{https://arxiv.org/abs/0906.2295}.
\bibitem{jones} J. P. Jones, \emph{Universal Diophantine equation},
Journal of Symbolic Logic \textbf{47}(3) (1982), 549--571.
\url{https://doi.org/10.2307/2273588}.
\end{thebibliography}
\endgroup
